\begin{table}[htbp]
\caption{Top 10 Brazilian coastal municipalities by the integrated compound-risk index. Ranks 1--3 are stable across a bootstrap over the 33 years of record, but the 90\textbackslash\{\}\% intervals of ranks 4--11 overlap, so the ordering within that band is not resolved; see the rank-uncertainty and aggregation-sensitivity tables (AUD-07 supplementary). Two entries carry a declared spatial-support caveat: Mag\textbackslash\{\}'e and Paraty sit inside sheltered bays and draw their hazard from open-shelf grid points 35 and 15 km away, so their hazard is imported from outside the embayment that shelters them; the recurrent flooding documented at both is fluvial and pluvial rather than wave-driven (AUD-04, AUD-05).}
\label{tab:top10-municipal-integrated-risk}
\begin{tabular}{rlllr}
\toprule
Rank & Municipality & State & UF & Risk index \\
\midrule
1 & São José do Norte & Rio Grande do Sul & RS & 0.566 \\
2 & Guaraqueçaba & Paraná & PR & 0.505 \\
3 & Magé & Rio de Janeiro & RJ & 0.468 \\
4 & Mangaratiba & Rio de Janeiro & RJ & 0.444 \\
5 & Paraty & Rio de Janeiro & RJ & 0.441 \\
6 & Guarujá & São Paulo & SP & 0.438 \\
7 & Balneário Gaivota & Santa Catarina & SC & 0.435 \\
8 & São Sebastião & São Paulo & SP & 0.432 \\
9 & Saquarema & Rio de Janeiro & RJ & 0.430 \\
10 & Passo de Torres & Santa Catarina & SC & 0.430 \\
\bottomrule
\end{tabular}
\end{table}
